% Transcribed from blueprint/src/parts/02-cited-facts.tex.
\section{What is cited}
\label{ch:cited}

This paper rests on five releases of the author's earlier work, cited at the printed number and
page. The line paper \progpaper{V-line v0.1} states the axioms on the line and classifies the
families that satisfy them. Two companion papers were released with it. The first, its
\emph{cone module} \progpaper{V-cone v0.1}, describes these families as a convex cone, with their
generators and their computation on the line. The second, its \emph{selection module}
\progpaper{V-selection v0.1}, treats locality and the non-creation of extrema. The isotropic
paper \progpaper{iso v0.1} carries the axioms and the classification to $\RR^d$. The affine paper
\progpaper{affine v0.1} supplies the frame families, which the supplement uses for the separability
of the classical kernels.

Section~\ref{sec:cited} states what this paper takes from these releases: the axioms, the two
classifications and the other statements that the proofs use.
Section~\ref{sec:cited-literature} states the facts of the literature: the \Levy--Khintchine
theory of infinitely divisible laws, Bernstein functions, \Polya{} frequency densities and one
integral of a Bessel function. Each fact of these theories that a proof takes as given is a
numbered theorem, its source and page
are named just before it, and a proof that uses it refers to it. Standard measure theory and
probability, such as \Levy{}'s continuity theorem, Steinhaus's theorem and the Radon--Nikodym
theorem, are used without a number.
Theorems~\ref{thm:cited-A30}, \ref{thm:cited-A13}, \ref{thm:cited-A19} and~\ref{thm:cited-A20}
orient the reader, and the remarks of Section~\ref{sec:line-theory} and the supplement compare
with the last three.

\subsection{The released modules: axioms, classifications and notation}
\label{sec:cited}

\subsubsection{The axioms on the line}
\label{subsec:cited-line}

A scale space keeps a record of a signal at every scale $t \ge 0$, the signal itself at $t = 0$,
and an operator $\Phi_{s,t}$ takes the record at scale $s$ to the record at the coarser scale $t$.
The cascade (A6) requires that two such steps compose to one. It replaces the semigroup law
$\Phi_{t+s} = \Phi_t\Phi_s$ of the classical theory, which ties the scale to a fixed additive
parametrization, and a family with (A6) is a \emph{hemigroup}. Covariance (A8) requires that a dilation
of the signal only relabel the scales, and (ND) excludes the identity.

% shared with blueprint def:cascade-measurement-family
\begin{definition}[symmetric cascade measurement families]
\label{def:cascade-measurement-family}

A \emph{symmetric cascade measurement family} is a family of linear operators
$\Phi_{s,t} : L^1(\RR) \to L^1(\RR)$, $0 \le s \le t < \infty$, such that
\begin{itemize}
\item[(A1)] each $\Phi_{s,t}$ is a bounded operator on $L^1(\RR)$;
\item[(A2)] $\Phi_{s,t}$ commutes with every translation $f \mapsto f(\cdot - a)$, $a \in \RR$;
\item[(A3)] $\Phi_{s,t}$ commutes with the reflection $f \mapsto f(-\,\cdot)$;
\item[(A4)] $\Phi_{s,t}$ maps nonnegative functions to nonnegative functions;
\item[(A5)] $\int\Phi_{s,t}f = \int f$ for every nonnegative $f \in L^1(\RR)$;
\item[(A6)] $\Phi_{t,t} = \Id$, and $\Phi_{s,t}\Phi_{r,s} = \Phi_{r,t}$ for all $r \le s \le t$
(the cascade);
\item[(A7)] for every $f \in L^1(\RR)$ the map $(s,t) \mapsto \Phi_{s,t}f$ is continuous from
$\{0 \le s \le t\}$ to $L^1(\RR)$;
\item[(A8)] there is a family $(S_\lambda)_{\lambda > 0}$ of increasing bijections of $[0,\infty)$
with $\mathcal{D}_\lambda\Phi_{s,t} = \Phi_{S_\lambda s, S_\lambda t}\mathcal{D}_\lambda$ for all
$0 \le s \le t$ and all $\lambda > 0$, where $(\mathcal{D}_\lambda f)(x) := \lambda^{-1}f(x/\lambda)$
(scale covariance);
\item[(ND)] $\Phi_{s,t} \ne \Id$ whenever $s < t$ (nondegeneracy).
\end{itemize}
\end{definition}

The first consequence of the axioms is that each operator is a convolution
\prog{V-line v0.1}{Lemma~4.1, pp.~22--23}.

% shared with blueprint thm:cited-A1
\begin{theorem}[each operator is a convolution]
\label{thm:cited-A1}

Let $\Phi_{s,t}$, $0 \le s \le t$, be linear operators on $L^1(\RR)$ with (A1), (A2), (A4) and
(A5) of Definition~\ref{def:cascade-measurement-family}. Then for each pair $0 \le s \le t$ there
is a unique probability measure $\mu_{s,t}$ on $\RR$ with
\[
  \Phi_{s,t}f = \mu_{s,t} * f, \qquad f \in L^1(\RR);
\]
under (A3) it is symmetric, invariant under $x \mapsto -x$. Conversely, every symmetric
probability measure on $\RR$ defines, by convolution, an operator with (A1)--(A5).
\end{theorem}

\begin{proof}
Existence and uniqueness are \prog{V-line v0.1}{Lemma~4.1, p.~22}, with the proof printed there
and proved again in the formal development:
under (A1), (A2), (A4) and (A5) each $\Phi_{s,t}$ is the convolution with exactly one probability
measure $\mu_{s,t}$, the equation $\Phi_{s,t}f = \mu_{s,t} * f$ holding in $L^1(\RR)$, that is,
almost everywhere.

For the symmetry, write $f^-(x) := f(-x)$ and $\mu^-$ for the image of $\mu = \mu_{s,t}$
under $x \mapsto -x$. The change of variables $y \mapsto -y$ gives
$\mu^- * f = (\mu * f^-)^-$ for every $f \in L^1(\RR)$. By (A3),
$\Phi_{s,t}f^- = (\Phi_{s,t}f)^-$, so $\mu^- * f = (\Phi_{s,t}f^-)^- =
\Phi_{s,t}f = \mu * f$ for every $f$. A probability measure is determined by its convolution
operator on $L^1(\RR)$: with $g(x) := (2\pi)^{-1/2}e^{-x^2/2}$, the transform of $\mu * g$ is
$\widehat\mu\,\widehat g$, and $\widehat g$ has no zero. So $\mu^- = \mu$.

For the converse let $\mu$ be a symmetric probability measure on $\RR$. The map
$f \mapsto \mu * f$ is linear, and $\|\mu * f\|_1 \le \|f\|_1$ by Tonelli's theorem, so it is a
bounded operator on $L^1(\RR)$, which is (A1). Since
$(\mu * f(\cdot - a))(x) = \int f(x - a - y)\,\mu(dy) = (\mu * f)(x - a)$, it commutes with
translation, which is (A2). Since $\mu^- = \mu$, the identity of the previous paragraph reads
$\mu * f^- = (\mu * f)^-$, which is (A3). If $f \ge 0$ almost everywhere, then
$\mu * f \ge 0$ almost everywhere, which is (A4), and Fubini's theorem gives
$\int \mu * f = \mu(\RR)\int f = \int f$, which is (A5).

Every step is machine-checked, the released lemma's included, and no fact of the literature is
cited.
\end{proof}

Under (A1)--(A7) the transform of the kernel is $\widehat\mu_{s,t} = e^{-g_{s,t}}$ with
\[
  g_{s,t}(\omega) = a(s,t)\,\omega^2 + \int_{(0,\infty)}(1 - \cos\omega x)\,\varpi_{s,t}(dx),
  \quad a(s,t) \ge 0, \quad \int_{(0,\infty)}\min(1, x^2)\,\varpi_{s,t}(dx) < \infty
\]
\prog{V-line v0.1}{Thm.~5.2, p.~26}. So every \emph{step} $\mu_{s,t}$ is a
symmetric \emph{infinitely divisible} law: for every $n$ it is the $n$-fold convolution power of
some law. The exponent of a step has a Gaussian part
$a(s,t)\omega^2$ and a jump part, and $\varpi_{s,t}$, its \emph{\Levy{} measure} folded onto
$(0,\infty)$, counts the jumps by their length. The pair $(a, \varpi)$ of such an exponent is
determined by the function \prog{V-line v0.1}{Prop.~2.3(3), p.~9}. The proof reads the
\Levy--Khintchine representation, Theorem~\ref{thm:cited-A39} of
Section~\ref{sec:cited-literature}.

% shared with blueprint thm:cited-A9
\begin{theorem}[the folded \Levy--Khintchine pair is unique]
\label{thm:cited-A9}

Let $a, a' \ge 0$, and let $\nu$ and $\nu'$ be measures on $(0,\infty)$ with
$\int_{(0,\infty)}\min(1, x^2)\,\nu(dx) < \infty$ and
$\int_{(0,\infty)}\min(1, x^2)\,\nu'(dx) < \infty$. If
\[
  a\,\omega^2 + \int_{(0,\infty)}(1 - \cos\omega x)\,\nu(dx)
  = a'\,\omega^2 + \int_{(0,\infty)}(1 - \cos\omega x)\,\nu'(dx)
  \qquad\text{for every } \omega \in \RR,
\]
then $a = a'$ and $\nu = \nu'$.
\end{theorem}

\begin{proof}
For a pair $(a, \nu)$ as in the statement let $\tilde\nu$ be the symmetric measure on $\RR$ with
$\tilde\nu(\{0\}) = 0$ and $\tilde\nu(B) = \tfrac12\nu(B \cap (0,\infty)) +
\tfrac12\nu(-B \cap (0,\infty))$. Then $\int\min(x^2, 1)\,\tilde\nu(dx) =
\int_{(0,\infty)}\min(x^2, 1)\,\nu(dx) < \infty$, so $(2a, \tilde\nu, 0)$ is a generating triplet
on $\RR$. In its \Levy--Khintchine bracket the imaginary part of the integrand,
$\sin\omega x - \omega x\mathbf 1_{|x| \le 1}$, is odd and $\tilde\nu$-integrable, so it integrates
to $0$. The real part $\cos\omega x - 1$ is even, so its $\tilde\nu$-integral is
$-\int_{(0,\infty)}(1 - \cos\omega x)\,\nu(dx)$. With the Gaussian term
$-\tfrac12\cdot 2a\,\omega^2$ the bracket is $-\psi(\omega)$, where $\psi$ is the function of the
pair, $\psi(\omega) = a\omega^2 + \int_{(0,\infty)}(1 - \cos\omega x)\,\nu(dx)$.

By Theorem~\ref{thm:cited-A39}(iii) there is an infinitely divisible law $\mu$ on $\RR$ with the triplet
$(2a, \tilde\nu, 0)$, so $\widehat\mu = e^{-\psi}$. By hypothesis the pair $(a', \nu')$ has the
same function $\psi$. So by the first paragraph the triplet $(2a', \tilde\nu', 0)$ also represents
$\widehat\mu$ by the formula of Theorem~\ref{thm:cited-A39}(i). By Theorem~\ref{thm:cited-A39}(ii) the two triplets
are equal: $a = a'$ and $\tilde\nu = \tilde\nu'$. On $(0,\infty)$ the measure $\tilde\nu$ is
$\tfrac12\nu$, so $\nu = 2\tilde\nu|_{(0,\infty)} = 2\tilde\nu'|_{(0,\infty)} = \nu'$.

Every step is machine-checked, and the cited fact is \ledger{A39}(ii)--(iii), at $d = 1$.
\end{proof}

The full axioms produce a classical class of laws. By the cascade the kernel $\mu_{0,t}$ is the
kernel $\mu_{0,s}$ at a finer scale convolved with the step $\mu_{s,t}$. By covariance
$\mu_{0,s}$ is a dilate of $\mu_{0,t}$. So each kernel is a shrunken copy of itself convolved with a
law, which is the property the next definition names [\sscite{sato1999levy}, Def.~15.1, p.~90].

% shared with blueprint def:self-decomposable
\begin{definition}[self-decomposable and semi-selfdecomposable laws]
\label{def:self-decomposable}

Let $\mu$ be a probability measure on $\RR^d$. It is \emph{self-decomposable} if, for every
$b > 1$, there is a probability measure $\varrho_b$ on $\RR^d$ such that
\[
  \widehat\mu(z) = \widehat\mu(b^{-1}z)\,\widehat\varrho_b(z), \qquad z \in \RR^d.
\]
It is \emph{semi-selfdecomposable} if there are some $b > 1$ and some infinitely divisible
probability measure $\varrho_b$ with the same identity. Then $b$ is called a \emph{span} of $\mu$, and
$\varrho_b$ its \emph{residual}.
\end{definition}

With $c := b^{-1} \in (0,1)$ the identity reads
$\widehat\mu(\omega) = \widehat\mu(c\omega)\,\widehat\varrho_b(\omega)$: $\mu$ is the law of
$cX + Y$ with $X \sim \mu$ and $Y \sim \varrho_b$ independent. The exponents of the
symmetric self-decomposable laws on the line are those of the form
\begin{equation}
\label{eq:admissible-exponent}
  F(\omega) = a\,\omega^2 + \int_0^\infty(1 - \cos\omega x)\,\frac{k(x)}{x}\,dx
\end{equation}
with $a \ge 0$, $k : (0,\infty) \to [0,\infty)$ nonincreasing and
$\int_0^\infty\min(1, x^2)\,k(x)\,x^{-1}\,dx < \infty$ \prog{V-line v0.1}{Prop.~2.13, p.~13}, the
integrability being that of a \Levy{} measure. The classification of the line is the main theorem
of the line paper \prog{V-line v0.1}{Thm.~7.3, p.~37}.

% shared with blueprint thm:cited-A4
\begin{theorem}[the classification on the line]
\label{thm:cited-A4}

A family $(\Phi_{s,t})_{0 \le s \le t}$ of linear operators on $L^1(\RR)$ satisfies (A1)--(A8)
and (ND) of Definition~\ref{def:cascade-measurement-family} if and only if there are an increasing
bijection $\chi$ of $[0,\infty)$, the \emph{gauge}, and a function $F \not\equiv 0$ of the form
\[
  F(\omega) = a\,\omega^2 + \int_0^\infty(1 - \cos\omega x)\,\frac{k(x)}{x}\,dx,
\]
with $a \ge 0$, $k \ge 0$ nonincreasing and
$\int_0^\infty\min(1, x^2)\,k(x)\,x^{-1}\,dx < \infty$, such that $\Phi_{s,t}f = \mu_{s,t} * f$ with
\[
  \widehat\mu_{s,t}(\omega) = \exp\bigl(-[F(\chi(t)\omega) - F(\chi(s)\omega)]\bigr),
  \qquad 0 \le s \le t,\ \omega \in \RR.
\]
Conversely, for every such pair $(\chi, F)$ there are probability measures $\mu_{s,t}$ with
these transforms, and the operators $f \mapsto \mu_{s,t} * f$ satisfy (A1)--(A8) and (ND).
The pair $(\chi, F)$ is unique once $\chi(1) = 1$.
\end{theorem}

In the \emph{canonical gauge} $\chi = \mathrm{id}$ the kernel at scale $t$ has transform
$e^{-F(t\omega)}$, and the family is the \emph{similarity family} of $F$. The functions $F$ of the
form (\ref{eq:admissible-exponent}) are the \emph{admissible exponents}, and $(a, k)$ their
\emph{data}. They form a convex cone \prog{V-cone v0.1}{Lemma~2.7, p.~12}. The Gaussian is the
exponent with $k = 0$. The \Matern{} exponent $\gamma\log(1 + \omega^2)$ of index $\gamma$ has
$a = 0$ and $k(x) = 2\gamma e^{-x}$.

\subsubsection{The axioms in $\RR^d$}
\label{subsec:cited-d}

The isotropic paper carries the axioms from the line to $\RR^d$, with rotation invariance in place
of symmetry, and classifies the families that satisfy them. This subsection states both, since
Section~\ref{ch:pyramid} weakens the first and compares its strata with the second. Fix $d \ge 1$.
For a measure $\mu$ on $\RR^d$ write
$\widehat\mu(\omega) = \int e^{i\langle\omega, x\rangle}\,\mu(dx)$, and for a linear map $A$ of
$\RR^d$ write $A_*\mu$ for the image of $\mu$ under $x \mapsto Ax$.

% shared with blueprint def:isotropic-cascade-family
\begin{definition}[isotropic cascade families]
\label{def:isotropic-cascade-family}

Let $d \ge 1$. An \emph{isotropic cascade family in dimension $d$} is a family
$(\mu_{s,t})_{0 \le s \le t}$ of probability measures on $\RR^d$ such that
\begin{itemize}
\item[(I1)] $\mu_{t,t} = \delta_0$, and $\mu_{r,t} = \mu_{s,t} * \mu_{r,s}$ whenever
$r \le s \le t$ (cascade);
\item[(I2)] for every $\omega \in \RR^d$ the map $(s,t) \mapsto \widehat\mu_{s,t}(\omega)$ is
continuous on $\{0 \le s \le t\}$ (continuity);
\item[(I3)] $Q_*\mu_{s,t} = \mu_{s,t}$ for every $Q \in \Orth(d)$ and every pair (rotation
invariance);
\item[(I4)] there is a family $(S_\lambda)_{\lambda > 0}$ of increasing bijections of
$[0,\infty)$ with $(\lambda I)_*\mu_{s,t} = \mu_{S_\lambda s, S_\lambda t}$ for every $\lambda > 0$
and every pair (dilation covariance);
\item[(ND)] $\mu_{s,t} \ne \delta_0$ whenever $s < t$ (nondegeneracy).
\end{itemize}
\end{definition}

The axioms are stated on the kernels. The isotropic paper also states them on operators on
$L^1(\RR^d)$, as (A1)--(A8) and (ND) with the rotations in place of the reflection
\prog{iso v0.1}{Def.~3.1, pp.~10--11}. The two definitions describe the same families, by the
next theorem \prog{iso v0.1}{Lemma~3.3, pp.~11--12}.

% shared with blueprint thm:cited-A30
\begin{theorem}[operators and kernels in dimension $d$]
\label{thm:cited-A30}

Let $d \ge 1$, and for a probability measure $\mu$ on $\RR^d$ and $f \in L^1(\RR^d)$ write
$(\mu * f)(x) := \int f(x - y)\,\mu(dy)$.
\begin{enumerate}
\item Let $\Phi$ be a linear operator on $L^1(\RR^d)$ that is bounded, commutes with every
translation $f \mapsto f(\cdot - a)$, $a \in \RR^d$, maps nonnegative functions to nonnegative
functions, and has $\int\Phi f = \int f$ for every nonnegative $f$. Then there is a unique
probability measure $\mu$ on $\RR^d$ with $\Phi f = \mu * f$ for every $f \in L^1(\RR^d)$.
Conversely, for every probability measure $\mu$ on $\RR^d$ the operator $f \mapsto \mu * f$ has
these four properties.
\item Let $\mu_n$ and $\mu$ be probability measures on $\RR^d$. Then
$\widehat\mu_n(\xi) \to \widehat\mu(\xi)$ for every $\xi \in \RR^d$ if and only if
$\mu_n * f \to \mu * f$ in $L^1(\RR^d)$ for every $f \in L^1(\RR^d)$.
\end{enumerate}
\end{theorem}

So (I1) is the cascade (A6), (I2) the continuity (A7) and (I4) the covariance (A8), with the same
bijections. (I3) is the rotation invariance, which in $\RR^d$ takes the place of the reflection
(A3). Let $\mathcal{A}_d$ be the set of admissible $F$ for
which $\omega \mapsto F(|\omega|)$, $\omega \in \RR^d$, is the exponent of a self-decomposable law
on $\RR^d$. Then $\mathcal{A}_1$ is the whole cone of admissible exponents and
$\mathcal{A}_{d+1} \subseteq \mathcal{A}_d$ \prog{iso v0.1}{Prop.~2.6, p.~8}. The classification
in $\RR^d$ is the main theorem of the isotropic paper \prog{iso v0.1}{Thm.~3.6, p.~13}.

% shared with blueprint thm:cited-A31
\begin{theorem}[the classification in $\RR^d$]
\label{thm:cited-A31}

Let $d \ge 1$, and let $\mathcal{A}_d$ be the set of admissible exponents $F$ (of the form
(\ref{eq:admissible-exponent})) for which $\omega \mapsto F(|\omega|)$ is the exponent of a
self-decomposable law on $\RR^d$. A family $(\mu_{s,t})_{0 \le s \le t}$ of probability measures
on $\RR^d$ is an isotropic cascade family in dimension $d$ if and only if there are a gauge
$\chi$, an increasing bijection of $[0,\infty)$, and an $F \in \mathcal{A}_d$ with
$F \not\equiv 0$ such that
\[
  \widehat\mu_{s,t}(\omega) = \exp\bigl(-[F(\chi(t)|\omega|) - F(\chi(s)|\omega|)]\bigr),
  \qquad 0 \le s \le t,\ \omega \in \RR^d.
\]
The pair is unique against any second such pair $(\chi', F')$ with $F' \in \mathcal{A}_d$, up to
$(\chi, F) \mapsto (c\chi, F(c^{-1}\,\cdot))$ with $c > 0$, so it is unique once $\chi(1) = 1$.
\end{theorem}

Its proof reduces an isotropic family to its marginals on a line \prog{iso v0.1}{\S~3.2, p.~12},
a reduction that Section~\ref{sec:pyramid-reduction} takes up under one dilation. Every admissible
pair is realized \prog{iso v0.1}{Lem.~3.8, p.~14}.

% shared with blueprint thm:cited-A34
\begin{theorem}[every pair is realized]
\label{thm:cited-A34}

Let $d \ge 1$, let $\chi$ be a gauge and let $F \in \mathcal{A}_d$ with $F \not\equiv 0$. Then
there is an isotropic cascade family $(\mu_{s,t})_{0 \le s \le t}$ in dimension $d$ with
\[
  \widehat\mu_{s,t}(\omega) = \exp\bigl(-[F(\chi(t)|\omega|) - F(\chi(s)|\omega|)]\bigr),
  \qquad 0 \le s \le t,\ \omega \in \RR^d.
\]
\end{theorem}

\begin{proof}
Let $\mu$ be the self-decomposable law on $\RR^d$ with $\widehat\mu(\omega) = e^{-F(|\omega|)}$,
which $F \in \mathcal{A}_d$ provides. For $0 < c < 1$ it has a residual $\varrho_c$, a probability
measure with $\widehat\mu(\omega) = \widehat\mu(c\omega)\,\widehat\varrho_c(\omega)$; put
$\varrho_0 := \mu$ and $\varrho_1 := \delta_0$. Since $\widehat\mu$ has no zero and $F(0) = 0$,
\[
  \widehat\varrho_c(\omega) = \exp\bigl(-[F(|\omega|) - F(c|\omega|)]\bigr), \qquad 0 \le c \le 1 .
\]
A gauge has $\chi(0) = 0$ and $\chi(t) > 0$ for $t > 0$. For $0 \le s \le t$ with $t > 0$ let
$\mu_{s,t}$ be the image of $\varrho_{\chi(s)/\chi(t)}$ under $x \mapsto \chi(t)x$, and let
$\mu_{0,0} := \delta_0$. The transform of $\mu_{s,t}$ at $\omega$ is that of $\varrho_{\chi(s)/\chi(t)}$
at $\chi(t)\omega$, which is the display.

The display gives the axioms of Definition~\ref{def:isotropic-cascade-family}. (I1): the exponents
of $\mu_{s,t}$ and $\mu_{r,s}$ add up to that of $\mu_{r,t}$, and a law is determined by its
transform. (I2): $F$ and $\chi$ are continuous. (I3): the transform depends on $\omega$ only through
$|\omega|$. (I4): with $S_\lambda := \chi^{-1}(\lambda\chi(\cdot))$, an increasing bijection of
$[0,\infty)$, the transform of $(\lambda I)_*\mu_{s,t}$ at $\omega$ is the display at
$(S_\lambda s, S_\lambda t)$. (ND): if $\mu_{s,t} = \delta_0$ with $s < t$, then $F(cr) = F(r)$ for
every $r \ge 0$, with $c := \chi(s)/\chi(t) < 1$; iterating, $F(r) = F(c^nr) \to F(0) = 0$, against
$F \not\equiv 0$.

Every step is machine-checked, in the affine paper's formal development, with no cited fact.
\end{proof}

Section~\ref{ch:pyramid} keeps (I1), (I2) and (ND), treats the rotation invariance (I3) as an
optional property, and requires (I4) at one dilation (Definition~\ref{def:family-dimension-d}). The class $\mathcal{A}_d$
has a polar description, by the \Levy{} measures of its laws \prog{iso v0.1}{Prop.~3.11(1)--(2),
pp.~15--16}. In the next definition the Gaussian coefficient $a_0$ enters the exponent as
$\tfrac12 a_0|\xi|^2$, in Sato's normalization. At $d = 1$ it is twice the $a$ of
(\ref{eq:admissible-exponent}) (Theorem~\ref{thm:cited-A35}(2)).

% shared with blueprint def:isotropic-origin-data
\begin{definition}[isotropic origin data]
\label{def:isotropic-origin-data}

Let $d \ge 1$, write $\sigma_{d-1}$ for the surface measure of the unit sphere
$S^{d-1} \subset \RR^d$ and $|S^{d-1}| := \sigma_{d-1}(S^{d-1})$. An \emph{isotropic origin datum
in dimension $d$} is a pair $(a_0, k_d)$ with $a_0 \ge 0$ and $k_d : (0,\infty) \to [0,\infty)$
nonincreasing, subject to
\[
  \int_0^1 k_d(r)\,r\,dr < \infty \qquad\text{and}\qquad \int_1^\infty k_d(r)\,\frac{dr}{r} < \infty .
\]
Its \emph{\Levy{} measure} is the rotation-invariant measure $\nu_d$ on $\RR^d \setminus \{0\}$
with radial density $f_d(r) = k_d(r)\bigl(|S^{d-1}|\,r^d\bigr)^{-1}$, so that
\[
  \int_{\RR^d} g\,d\nu_d = \frac{1}{|S^{d-1}|}\int_{S^{d-1}}\sigma_{d-1}(d\eta)
  \int_0^\infty g(r\eta)\,k_d(r)\,\frac{dr}{r}
\]
for every nonnegative Borel $g$ vanishing at the origin. Its \emph{exponent} is
\[
  F_d(\xi) = \tfrac12 a_0|\xi|^2 + \int_{\RR^d}\bigl(1 - \cos\langle\xi, x\rangle\bigr)\,\nu_d(dx),
  \qquad \xi \in \RR^d,
\]
and its \emph{law} is the symmetric infinitely divisible law $\mu$ on $\RR^d$ with
$\widehat\mu = e^{-F_d}$. The function $k_d$ is the datum's \emph{polar profile}.
\end{definition}

% shared with blueprint thm:cited-A35
\begin{theorem}[the polar form of $\mathcal{A}_d$]
\label{thm:cited-A35}

Let $d \ge 1$ and $F : [0,\infty) \to [0,\infty)$.
\begin{enumerate}
\item $\xi \mapsto F(|\xi|)$ is the exponent of a self-decomposable law on $\RR^d$ if and only if
it is the exponent $F_d$ of an isotropic origin datum $(a_0, k_d)$ in dimension $d$
(Definition~\ref{def:isotropic-origin-data}):
\[
  F(|\xi|) = \tfrac12 a_0|\xi|^2 + \int_{\RR^d}\bigl(1 - \cos\langle\xi, x\rangle\bigr)\,
  \frac{k_d(|x|)}{|S^{d-1}|\,|x|^d}\,dx, \qquad \xi \in \RR^d.
\]
The datum is then unique, $k_d$ up to modification on a Lebesgue-null set.
\item If this holds, $F$ is admissible, of the form (\ref{eq:admissible-exponent}) with Gaussian
coefficient $a = \tfrac12 a_0'$ and profile $k = k_1$ for its datum $(a_0', k_1)$ in dimension $1$. So
$\mathcal{A}_d$ is the set of $F$ for which $F(|\cdot|)$ is the exponent of an isotropic origin
datum in dimension $d$.
\end{enumerate}
\end{theorem}

\subsubsection{The generator form}
\label{subsec:cited-generator}

The scale space of an admissible exponent solves an evolution equation in the scale. Let $F$ be
admissible with data $(a, k)$, put $\kappa^{\mathrm C} := -dk$, a measure on $(0,\infty)$, and let
$u(t,\cdot) = \mu_{0,t} * f$ be the scale space of a Schwartz function $f \in \mathcal{S}(\RR)$ in
the canonical gauge. Then $\partial_t u = \mathcal{H}_t u$ with
\begin{equation}
\label{eq:generator-form}
  \mathcal{H}_t g(x) = 2at\,g''(x) + \frac1t\int_{(0,\infty)}\delta^2_{tv}g(x)\,
  \kappa^{\mathrm C}(dv), \qquad
  \delta^2_h g(x) := \tfrac12\bigl(g(x - h) + g(x + h)\bigr) - g(x)
\end{equation}
\prog{V-cone v0.1}{Prop.~6.2, display (6.2), p.~56}. Write
$\Cin(z) := \int_0^z(1 - \cos v)\,dv/v$, the entire cosine integral. The \emph{Cin ray of step
$\tau$} is the exponent $\Cin(\tau\omega)$, which is (\ref{eq:admissible-exponent}) with $a = 0$ and
the profile $k = \mathbf{1}_{(0,\tau)}$ \prog{V-cone v0.1}{Lemma~3.1(1), p.~13}. Its generator is a
single symmetric second difference \prog{V-cone v0.1}{Prop.~6.4(5), p.~58}.

% shared with blueprint thm:cited-A13
\begin{theorem}[the generator of a Cin ray]
\label{thm:cited-A13}

Let $\tau > 0$, and let $F$ be the admissible exponent with data $(0, \mathbf{1}_{(0,\tau)})$, the
Cin ray of step $\tau$. In the canonical gauge its generator is a single symmetric second
difference: for $f \in \mathcal{S}(\RR)$ the scale space $u(t,\cdot) = \mu_{0,t} * f$ satisfies,
for every $t > 0$,
\[
  \partial_t u(t,\cdot) = \frac1t\,\delta^2_{t\tau}u(t,\cdot), \qquad
  \delta^2_h g(x) = \tfrac12\bigl(g(x - h) + g(x + h)\bigr) - g(x).
\]
\end{theorem}

So (\ref{eq:generator-form}) writes every admissible evolution as a Gaussian term, with the weight
$a$, plus a positive superposition of these \emph{atomic generators}, with the weight
$\kappa^{\mathrm C}$. The same weights, with $\kappa^{\mathrm C}$ written $\varpi$ below, give $F$
itself, and the Gaussian ray and the Cin rays are the extreme rays of the cone of admissible
exponents \prog{V-cone v0.1}{Prop.~3.3, pp.~15--16}.

% shared with blueprint thm:cited-A15
\begin{theorem}[the Choquet structure of the admissible cone]
\label{thm:cited-A15}

The superposition map $(a, \varpi) \mapsto F$,
\[
  F(\omega) = a\omega^2 + \int_{(0,\infty)}\Cin(\tau\omega)\,\varpi(d\tau), \qquad \omega \in \RR,
\]
is a bijection from the pairs $(a, \varpi)$, $a \ge 0$ and $\varpi$ a positive Borel measure on
$(0,\infty)$ with $\int\min\bigl(\tau^2, 1 + \log_+\tau\bigr)\,\varpi(d\tau) < \infty$, onto the
admissible exponents. Its inverse takes the exponent with data $(a, k)$ to $(a, \varpi)$ with
$\varpi = -dk$, that is $k(x) = \varpi((x,\infty))$ for almost every $x > 0$. Consequently the
extreme rays of the cone of admissible exponents are the Gaussian ray $\omega \mapsto \omega^2$
and the Cin rays $\omega \mapsto \Cin(\tau\omega)$, $\tau > 0$, and every admissible exponent is a
unique superposition of Cin rays and the Gaussian term $a\omega^2$.
\end{theorem}

\subsubsection{Other statements of the released modules}
\label{subsec:cited-released}

The paper and the supplement also cite the following statements of the released modules. The list
is an index: each is restated, with what it says, where it is used.
\begin{itemize}
\item \prog{V-line v0.1}{Remark~6.7, pp.~32--33}: the gauge freedom of the scale coordinate.
\item \prog{V-cone v0.1}{Prop.~3.3, pp.~15--16}: Theorem~\ref{thm:cited-A15}.
\item \prog{V-cone v0.1}{\S~4, p.~24}: the subordination bridge, which reads a family as Brownian
motion run at a random time.
\item \prog{V-cone v0.1}{Prop.~5.7, p.~48}: the symmetric Thorin subclass.
\item \prog{V-cone v0.1}{Prop.~6.4(4)--(5), p.~58}: the generator of a Thorin family as a mixture of
differences of Laplace kernels, and that of a Cin ray as one symmetric second difference.
\item \prog{V-selection v0.1}{Thm.~4.2, pp.~30--31}: non-enhancement at global extrema for every family,
at local extrema on the Gaussian ray only; and its Cor.~4.4, p.~33: off the Gaussian ray the scale
space enhances a local maximum.
\item \prog{V-selection v0.1}{\S~4.3, pp.~34--35, and Prop.~4.5, p.~35}: the \Polya{} frequency
members and their counts, restated below as Theorems~\ref{thm:cited-A22} and~\ref{thm:cited-A23}.
\item \prog{V-selection v0.1}{Thm.~4.10, p.~43}: scale-monotone non-creation characterizes the
Gaussian.
\item \prog{V-selection v0.1}{Remark~4.15, pp.~50--51}: four readings of non-creation, and the
nested classes of families they select.
\item \prog{V-selection v0.1}{Prop.~5.6, p.~60}: the Helmholtz identity.
\item \prog{V-selection v0.1}{\S~6, p.~69}: diffusion creates extrema in the plane.
\item \prog{iso v0.1}{Prop.~3.12, pp.~16--17}: the subordination bridge in $\RR^d$; the
\Matern{}, stable and Student-$t$ kernels.
\item \prog{iso v0.1}{Prop.~3.22, p.~27}: the polar profile one dimension down;
$\mathcal{A}_{d+1} \subsetneq \mathcal{A}_d$ for every $d$.
\item \prog{affine v0.1}{Prop.~F.1, p.~93}: the frame families.
\end{itemize}
The cone module computes the \Matern{} family on the line by forward--backward first-order sections
in its Proposition~7.2, pp.~60--61, and interpolates the pyramid of a rational step in its
Proposition~7.5(2), p.~65.

% shared with blueprint thm:cited-A19
\begin{theorem}[the \Matern{} cascade by first-order sections]
\label{thm:cited-A19}

Let $\gamma \ge 1$ be an integer, and consider the \Matern{} family on the line in the canonical
gauge, $e^{-F(t\omega)} = (1 + t^2\omega^2)^{-\gamma}$. For $t > 0$ let $e^\pm_t$ be
the one-sided exponential kernels $t^{-1}e^{-|x|/t}$ carried by $\pm(0,\infty)$, and
$L_t$ the Laplace kernel $(2t)^{-1}e^{-|x|/t}$. Let $0 < t_k < t_{k+1}$ be adjacent knots
and $r := t_k/t_{k+1} \in (0,1)$. Then the increment has transform
\[
  \widehat\mu_{t_k, t_{k+1}}(\omega) = \Bigl(\frac{1 + t_k^2\omega^2}{1 + t_{k+1}^2\omega^2}\Bigr)^{\gamma},
\]
and it is a cascade of $\gamma$ identical causal first-order sections and $\gamma$ identical
anticausal ones. Each causal section is the convolution with
$r\,\delta_0 + (1 - r)\,e^+_{t_{k+1}}$, a mixture of no displacement and a one-sided
exponential displacement, with the one-state realization
\[
  w' = \frac{y - w}{t_{k+1}}, \qquad v = r\,y + (1 - r)\,w,
\]
run forward in $x$ (input $y$, state $w$, output $v$, the prime the derivative in $x$), the
solution $w$ being the one in $L^1$, with $w(-\infty) = 0$. Each anticausal
section is its mirror image, run backward. For $\gamma = 1$ the two-sided increment is
\[
  r^2\delta_0 + r(1 - r)\bigl(e^+_{t_{k+1}} + e^-_{t_{k+1}}\bigr)
  + (1 - r)^2\,L_{t_{k+1}} .
\]
\end{theorem}

% shared with blueprint thm:cited-A20
\begin{theorem}[rational increments and the pyramid they interpolate]
\label{thm:cited-A20}

Let $R$ be a symmetric rational function that is the transform of an infinitely divisible law on
$\RR$, with folded \Levy{} profile $k_R$. Then for each $q > 1$ there is exactly one function $F$
that is continuous with $F(0) = 0$ and satisfies $F(q\omega) - F(\omega) = -\log R(\omega)$,
namely
\[
  e^{-F(\omega)} = \prod_{m \ge 1}R(q^{-m}\omega), \qquad\text{with folded profile}\qquad
  k(x) = \sum_{m \ge 1}k_R(q^m x);
\]
it is admissible exactly when this $k$ has a nonincreasing version ($k_R$, and with it $k$, being
fixed only up to a null set), in particular whenever $k_R$ is nonincreasing. The
\Matern{} pyramid is the case $R(\omega) = (1 + \omega^2)^{-1}$.
\end{theorem}

A \emph{\Polya{} frequency density} is a probability density $\phi$ on $\RR$ whose translation
kernel $(x, y) \mapsto \phi(x - y)$ is totally positive: every determinant
$\det[\phi(x_i - y_j)]_{i,j \le n}$ with $x_1 < \dots < x_n$ and $y_1 < \dots < y_n$ is
nonnegative [\sscite{karlin1968total}, Ch.~7, \S 1, p.~332]. Convolution with such a density is
\emph{variation-diminishing}: it does not increase the number of sign changes of a function. This
is the property that the counting statements use. The selection module shows that the families of
the exponents below have \Polya{} frequency kernels, and counts what they create from a signal, in
its Proposition~4.5(1)--(2), p.~35. Both facts of the next theorem are classical. A density with
the transform $e^{-F}$, for $F$ of the form below, is a \Polya{} frequency density by Schoenberg's
theorem [\sscite{karlin1968total}, Ch.~7, (1.6), p.~335, and Thm.~3.2(a), p.~345], with the numbers
$\pm1/\theta_i$ in the product. A \Polya{} frequency kernel is variation-diminishing
[\sscite{karlin1968total}, Ch.~5, Thm.~3.1, p.~233]. Karlin states it for bounded Borel inputs,
with the input's sign changes counted in the essential sense and the output's pointwise. The
selection module extends it to integrable inputs and to the essential count for the output
\prog{V-selection v0.1}{Prop.~4.5, p.~35}. The
theorem restates the selection module's proposition, which is the statement the formal development
takes as an axiom.

% shared with blueprint thm:cited-A22
\begin{theorem}[the \Polya{} frequency members]
\label{thm:cited-A22}

For a function $g$ on $\RR$ write $S^-(g)$ for the number of sign changes of $g$, in the essential
sense.
Let $F \not\equiv 0$ be an admissible exponent of the form
\[
  F(\omega) = a\,\omega^2 + \sum_i\log\Bigl(1 + \frac{\omega^2}{\theta_i^2}\Bigr), \qquad
  a \ge 0,\quad \theta_i > 0,\quad \sum_i\theta_i^{-2} < \infty,
\]
over countably many $\theta_i$. Then the similarity family of $F$, with kernels
$\widehat\mu_{0,t} = e^{-F(t\,\cdot)}$, has \Polya{} frequency kernels, so that convolution with
$\mu_{0,t}$ is variation-diminishing: $S^-(\mu_{0,t} * g) \le S^-(g)$ for every $t > 0$ and every
integrable or bounded $g$.
\end{theorem}

% shared with blueprint thm:cited-A23
\begin{theorem}[the counting clause]
\label{thm:cited-A23}

For a family as in Theorem~\ref{thm:cited-A22} and $f$ the primitive of an integrable $f'$, the
scale space $u(t,\cdot) = \mu_{0,t} * f$ has, at every scale $t > 0$, no more crossings of any
value $c$ than the signal and no more changes of monotonicity:
\[
  S^-\bigl(u(t,\cdot) - c\bigr) \le S^-(f - c) \quad (c \in \RR), \qquad
  S^-\bigl(\partial_x u(t,\cdot)\bigr) \le S^-(f').
\]
\end{theorem}

\begin{proof}
The primitive $f = C + \int_0^x f'$ is bounded by $|C| + \|f'\|_1$ and continuous, so
$f - c$ is a bounded measurable test function and $\mu_{0,t} * f - c = \mu_{0,t} * (f - c)$ at
every point. By Fubini's theorem on $[0, x] \times \mu_{0,t}$,
$u(t, x) = u(t, 0) + \int_0^x (\mu_{0,t} * f')(z)\,dz$ at every $x$, with $\mu_{0,t} * f'$ integrable;
by the Lebesgue differentiation theorem $\partial_x u(t,\cdot) = \mu_{0,t} * f'$ almost everywhere,
and $S^-$, a count in the essential sense, does not see the null set. The variation-diminishing
clause of Theorem~\ref{thm:cited-A22}, applied to the bounded $f - c$ and to the integrable $f'$,
gives the two inequalities.
\end{proof}

\subsubsection{Notation}
\label{subsec:cited-notation}

This paper writes the law with transform $(1 + r^2|\omega|^2)^{-\gamma}$ as the \Matern{} law
\emph{of index $\gamma$ at range $r$}. The isotropic paper's \Matern{} kernel of index $\gamma$ at
scale $t$ has transform $(1 + t^2|\omega|^2/2)^{-\gamma}$ \prog{iso v0.1}{Prop.~3.12(2), p.~16},
so it is the law at range $t/\sqrt2$. The cone module's \Matern{} kernel at scale $t$ has transform
$(1 + t^2\omega^2)^{-\gamma}$ \prog{V-cone v0.1}{Prop.~5.2(2), p.~43}, the law at range $t$.

The transform $\widehat\mu$ and the image $A_*\mu$ are those of Section~\ref{subsec:cited-d}. The
dilation of measures is $D_c := (cI)_*$ for $c > 0$, so that
$\widehat{D_c\mu}(\omega) = \widehat\mu(c\omega)$. On the line it is written $T_c$, as in
[\sscite{sato1999levy}, (13.8), p.~75]: $(T_c\mu)(B) = \mu(c^{-1}B)$ for Borel sets $B \subseteq \RR$. The same
notation is used for a measure on $(0,\infty)$, such as a folded \Levy{} measure, since
$x \mapsto cx$ maps $(0,\infty)$ onto itself. The dilation of functions $\mathcal{D}_\lambda$ is the
one of (A8), mass-preserving. The ratio of a pyramid is $q > 1$, and $\ell := \log q$.

The Gaussian part has two normalizations. The paper's \Levy{} data of a
law on $\RR^d$ are a Gaussian part $A$ and a \Levy{} measure $\nu$
(Lemma~\ref{lem:symmetric-levy-khintchine}). The Gaussian part enters the exponent as
$\langle\omega, A\omega\rangle$, so at $d = 1$ it is the $a$ of (\ref{eq:admissible-exponent}). Sato's Gaussian covariance matrix, which enters as
$\tfrac12\langle z, A^{\mathrm S}z\rangle$ in the theorems of Section~\ref{subsec:cited-sato}, is
written $A^{\mathrm S}$, so $A^{\mathrm S} = 2A$. The Gaussian coefficient $a_0$ of an isotropic
origin datum (Definition~\ref{def:isotropic-origin-data}) is in Sato's normalization,
$A^{\mathrm S} = a_0I$. On the line the folded \Levy{} measure is written $\varpi$
and the measure of (\ref{eq:generator-form}) $\kappa^{\mathrm C}$, the two related by
$\varpi(dx) = k(x)\,dx/x$ with $k(x) = \kappa^{\mathrm C}((x,\infty))$. The generator of a family
is $\mathcal{H}_t$, and $\mathcal{A}_d$ is the class of exponents admitted in dimension $d$. The
knots, the knot kernels and the ladder of a pyramid are defined after
Lemma~\ref{lem:family-dilations}.

The \Levy{} data have a reading as operators. A symmetric infinitely divisible law with exponent
$\Psi(\omega) = \langle\omega, A\omega\rangle + \int(1 - \cos\langle\omega, x\rangle)\,\nu(dx)$ is
the time-one law of the convolution semigroup with generator
\[
  g \mapsto \nabla\cdot(A\nabla g) + \int \Big(\tfrac12\big(g(\cdot + x) + g(\cdot - x)\big) - g\Big)\,\nu(dx).
\]
The Gaussian part is a diffusion. The \Levy{} measure weights the symmetric second differences of
span $x$, so $\nu$ says how much smoothing acts at each displacement. The probabilistic word for this
is a jump of size $x$, and the paper uses it.

Two families carry the name \Matern{}. The law $L$ is the kernel of the one and the octave step of
the other. The \Matern{} family of index $1$ has the kernel $L_t$ at the scale $t$, the \Matern{}
law of index $1$ at range $t$. Its octave step has the transform
$(1 + t^2|\omega|^2)/(1 + q^2t^2|\omega|^2)$ (Proposition~\ref{prop:principle-other-readings}(1)),
so it is the mixture $q^{-2}\delta_0 + (1 - q^{-2})L_{qt}$, by the identity of
Proposition~\ref{prop:recursion-helmholtz}(1), which holds in every dimension.
The \Matern{} pyramid of Section~\ref{ch:matern} has the octave step $L_{ct}$, one Bessel potential,
and its kernel at the scale $t$ is the convolution of the steps $L_{ctq^{-i}}$, $i \ge 1$, below
it. Both families are covariant under every dilation.

One example puts the normalizations side by side. On the line, take the Gaussian law of variance
$s^2$ convolved with the compound Poisson law of jumps $\pm x$, each at rate $\lambda$. Its exponent
is $\tfrac12 s^2\omega^2 + 2\lambda(1 - \cos\omega x)$. The paper's data are $A = a = s^2/2$ and
$\nu = \lambda(\delta_x + \delta_{-x})$. Sato's are $A^{\mathrm S} = s^2$ and the same $\nu$. The
folded \Levy{} measure is $\varpi = 2\lambda\delta_x$. In $\RR^d$ the radial counterpart has the
Gaussian part $\tfrac12 s^2 I$ and $\nu = 2\lambda\sigma_x$, with $\sigma_x$ the uniform law on the
sphere of radius $x$, which is the line's $\nu$ at $d = 1$. Its exponent is
$\tfrac12 s^2|\omega|^2 + 2\lambda(1 - \Omega_d(x|\omega|))$, with
$\Omega_d(r) := \int\cos(r\eta_1)\,\sigma_1(d\eta)$. An isotropic origin datum with this Gaussian part
has $a_0 = s^2$. Its \Levy{} measure has a density, which $2\lambda\sigma_x$ lacks.

Some letters carry more than one meaning, each fixed where it is introduced.

\begin{center}
\small
\begin{tabular}{>{\raggedright\arraybackslash}p{0.17\textwidth}>{\raggedright\arraybackslash}p{0.74\textwidth}}
\toprule
Symbol & Meanings \\
\midrule
$D_c$ & the dilation of measures, written $T_c$ on the line \\
$\mathcal{D}_\lambda$ & the dilation of functions, mass-preserving (A8) \\
$D$, $E$ & balls in the statements of Section~\ref{subsec:cited-sato} \\
$D^k$ & a derivative of order $k$ (Section~\ref{ch:matern}) \\
$D_h$, $D_5$, $D_9$ & the symbols of the discrete Laplacians, not dilations
(Section~\ref{sec:discrete-helmholtz}) \\
$L_r$ & the \Matern{} law of index $1$ at range $r$; the operator $f \mapsto L_r * f$; on the line
also its density \\
$\sigma_{d-1}$ & the surface measure of the unit sphere, not normalized (Definition~\ref{def:isotropic-origin-data}) \\
$\sigma_r$ & the uniform probability law on the sphere of radius $r$ (Section~\ref{ch:pyramid});
at $d = 2$, $\sigma_1$ is this law at $r = 1$, not the surface measure \\
$\sigma$ & the step of a random walk on a lattice; the symbol $\sigma_t$ of the generator
(Proposition~\ref{prop:matern-generator-octave}), which the supplement writes $h_t$ on the line
(Theorem~\ref{thm:matern-pyramid-line}(3)); in the supplement, a measure on a cylinder \\
$\sigma_n$ & coefficients in the proof of Proposition~\ref{prop:sampled-laplace-exact} \\
$\mu$ & with indices, a kernel or a step; alone, in Sections~S2.2 and~S2.3 of the supplement, a
step measure \\
\bottomrule
\end{tabular}
\end{center}

\subsection{The facts of the literature}
\label{sec:cited-literature}

The proofs read facts from three sources. They are the \Levy--Khintchine theory of infinitely
divisible laws, in the book of Sato [\sscite{sato1999levy}], Schoenberg's representation of
\Polya{} frequency densities, in Karlin's book [\sscite{karlin1968total}], and one integral
representation of the modified Bessel function $K_\nu$ [\sscite{dlmf2026}]. A fourth source,
the book of Schilling, Song and Vondra\v{c}ek on Bernstein functions
[\sscite{schilling2012bernstein}], supplies the terms in which Section~\ref{ch:matern} describes the
steps of the \Matern{} pyramid. Each fact is stated at its source's letter, or narrower where the
proofs use less, and the sentence before it says where. Theorem~\ref{thm:cited-A50} is proved in
the paragraph before it.

\subsubsection{Infinitely divisible laws}
\label{subsec:cited-sato}

A probability measure $\mu$ on $\RR^d$ is infinitely divisible, as on the line, if for every $n$ it
is the $n$-fold convolution power of some probability measure. Write $D := \{x \in \RR^d : |x| \le 1\}$
for the closed unit ball. Sato's Gaussian covariance matrix is written $A^{\mathrm S}$, twice the
paper's Gaussian part (Section~\ref{subsec:cited-notation}). The first theorem is the
\Levy--Khintchine representation [\sscite{sato1999levy}, Thm.~8.1(i)--(iii), pp.~37--38], with the
names of [\sscite{sato1999levy}, Def.~8.2, p.~38].

% shared with blueprint thm:cited-A39
\begin{theorem}[the \Levy--Khintchine representation]
\label{thm:cited-A39}

\begin{enumerate}
\item[(i)] If $\mu$ is an infinitely divisible probability measure on $\RR^d$, then
\[
  \widehat\mu(z) = \exp\Bigl[-\tfrac12\langle z, A^{\mathrm S}z\rangle + i\langle\gamma, z\rangle
  + \int_{\RR^d}\bigl(e^{i\langle z, x\rangle} - 1 - i\langle z, x\rangle\mathbf{1}_D(x)\bigr)\,
  \nu(dx)\Bigr], \qquad z \in \RR^d,
\]
where $A^{\mathrm S}$ is a symmetric nonnegative-definite $d \times d$ matrix, $\gamma \in \RR^d$, and $\nu$
is a measure on $\RR^d$ with $\nu(\{0\}) = 0$ and $\int_{\RR^d}\min(|x|^2, 1)\,\nu(dx) < \infty$.
\item[(ii)] The representation of $\widehat\mu$ in (i) by $A^{\mathrm S}$, $\nu$ and $\gamma$ is unique.
\item[(iii)] Conversely, for every symmetric nonnegative-definite $d \times d$ matrix $A^{\mathrm S}$, every
measure $\nu$ with $\nu(\{0\}) = 0$ and $\int\min(|x|^2, 1)\,\nu(dx) < \infty$, and every
$\gamma \in \RR^d$, there is an infinitely divisible probability measure $\mu$ whose transform is
given by the formula of (i).
\end{enumerate}
The triple $(A^{\mathrm S}, \nu, \gamma)$ is the \emph{generating triplet} of $\mu$, $A^{\mathrm S}$ its
\emph{Gaussian covariance matrix} and $\nu$ its \emph{\Levy{} measure}.
\end{theorem}

The transform of an infinitely divisible law has no zero [\sscite{sato1999levy}, Lemma~7.5,
p.~32].

% shared with blueprint thm:cited-A40
\begin{theorem}[the transform has no zero]
\label{thm:cited-A40}

If $\mu$ is an infinitely divisible probability measure on $\RR^d$, then $\widehat\mu(z) \ne 0$ for
every $z \in \RR^d$.
\end{theorem}

Limits of row sums of a null array are infinitely divisible, by Khintchine's theorem
[\sscite{sato1999levy}, Thm.~9.3, p.~47], with the null array of Definition~9.2 on the same page.

% shared with blueprint thm:cited-A41
\begin{theorem}[Khintchine's theorem]
\label{thm:cited-A41}

Let $\{Z_{nk} : k = 1, \dots, r_n;\ n = 1, 2, \dots\}$ be random variables on $\RR^d$ such that, for
each $n$, $Z_{n1}, \dots, Z_{nr_n}$ are independent, and, for every $\varepsilon > 0$,
\[
  \lim_{n \to \infty}\ \max_{1 \le k \le r_n}\Pr\bigl[|Z_{nk}| > \varepsilon\bigr] = 0
\]
(a \emph{null array}), and let $S_n := \sum_{k=1}^{r_n}Z_{nk}$ be its row sums. If, for some
$b_n \in \RR^d$, the law of $S_n - b_n$ converges to a probability measure $\mu$, then $\mu$ is
infinitely divisible.
\end{theorem}

The triplet of a linear image is [\sscite{sato1999levy}, Prop.~11.10, p.~65], printed there for a
\Levy{} process and read here for its law at time $1$. It follows from the form of the transform in
Theorem~\ref{thm:cited-A39}(i) by the change of variables $y = Ux$.

% shared with blueprint thm:cited-A42
\begin{theorem}[the linear image of an infinitely divisible law]
\label{thm:cited-A42}

Let $\mu$ be an infinitely divisible probability measure on $\RR^d$ with generating triplet
$(A^{\mathrm S}, \nu, \gamma)$, and let $U$ be an $n \times d$ matrix. Then the image $U_*\mu$ of $\mu$ under
$x \mapsto Ux$ is infinitely divisible on $\RR^n$, with generating triplet
\[
  A^{\mathrm S}_U = UA^{\mathrm S}U^{\tp}, \qquad \nu_U = \bigl[\nu U^{-1}\bigr]_{\RR^n \setminus \{0\}}, \qquad
  \gamma_U = U\gamma + \int Ux\,\bigl(\mathbf{1}_E(Ux) - \mathbf{1}_D(x)\bigr)\,\nu(dx),
\]
where $(\nu U^{-1})(B) = \nu(\{x : Ux \in B\})$, $[\,\cdot\,]_{\RR^n \setminus \{0\}}$ is the
restriction to $\RR^n \setminus \{0\}$, and $D$ and $E$ are the closed unit balls of $\RR^d$ and
$\RR^n$.
\end{theorem}

\begin{proof}
Throughout, $\langle \xi, Ux\rangle = \langle U^{\tp}\xi, x\rangle$ for $x \in \RR^d$ and
$\xi \in \RR^n$, and $K$ is the operator norm of $U$, so that $|Ux| \le K|x|$. Write
$\psi_{A^{\mathrm S},\nu,\gamma}(z)$ for the bracket of Theorem~\ref{thm:cited-A39}(i), and
$g_z(x) = e^{i\langle z,x\rangle} - 1 - i\langle z,x\rangle\mathbf{1}_D(x)$ for its integrand on
$\RR^d$, with $\mathbf{1}_E$ in place of $\mathbf{1}_D$ on $\RR^n$.

\emph{Infinite divisibility.} Let $k \ge 1$ and let $\mu_k$ be a probability measure with
$\mu = \mu_k^{*k}$. The transform of $(U_*\mu_k)^{*k}$ at $\xi$ is
$\widehat{\mu_k}(U^{\tp}\xi)^k = \widehat\mu(U^{\tp}\xi)$, which is the transform of $U_*\mu$ at $\xi$; so
$U_*\mu = (U_*\mu_k)^{*k}$, and $U_*\mu$ is infinitely divisible.

\emph{The correction is integrable.} Put $c(x) = Ux\,(\mathbf{1}_E(Ux) - \mathbf{1}_D(x))$. It
vanishes unless exactly one of $|x| \le 1$ and $|Ux| \le 1$ holds. If $|x| > 1 \ge |Ux|$, then
$|c(x)| = |Ux| \le 1 = \min(|x|^2, 1)$. If $|x| \le 1 < |Ux|$, then $1 < |Ux| \le K|x|$, so
$|c(x)| \le K|x| \le (K|x|)^3 \le K^3|x|^2 = K^3\min(|x|^2, 1)$. Hence
$|c(x)| \le (K^3 + 1)\min(|x|^2, 1)$ for every $x$, and $c$ is $\nu$-integrable.

\emph{The triplet.} The matrix $UA^{\mathrm S}U^{\tp}$ is symmetric. It is nonnegative-definite, as
$\langle \xi, UA^{\mathrm S}U^{\tp}\xi\rangle = \langle U^{\tp}\xi, A^{\mathrm S}U^{\tp}\xi\rangle$. The measure
$\nu_U$ has no mass at $0$, and since $\min(|Ux|^2, 1) \le (K^2 + 1)\min(|x|^2, 1)$,
\[
  \int \min(|y|^2, 1)\,\nu_U(dy) \le \int \min(|Ux|^2, 1)\,\nu(dx)
  \le (K^2 + 1)\int \min(|x|^2, 1)\,\nu(dx) < \infty .
\]
So $(A^{\mathrm S}_U, \nu_U, \gamma_U)$ is a generating triplet on $\RR^n$ in the sense of
Theorem~\ref{thm:cited-A39}(i).

\emph{The transform.} Fix $\xi \in \RR^n$. The integrand $e^{i\langle \xi,y\rangle} - 1 - i\langle \xi,y\rangle\mathbf{1}_E(y)$
vanishes at $y = 0$, so its integral against $\nu_U$ is its integral against the image of $\nu$,
which is $\int (e^{i\langle \xi,Ux\rangle} - 1 - i\langle \xi,Ux\rangle\mathbf{1}_E(Ux))\,\nu(dx)$. Since
$\langle \xi, Ux\rangle = \langle U^{\tp}\xi, x\rangle$, the integrand here is
$g_{U^{\tp}\xi}(x) - i\langle \xi, c(x)\rangle$, a difference of two $\nu$-integrable functions, so the
integral is $\int g_{U^{\tp}\xi}\,d\nu - i\langle \xi, \int c\,d\nu\rangle$. With
$\langle \xi, A^{\mathrm S}_U \xi\rangle = \langle U^{\tp}\xi, A^{\mathrm S}U^{\tp}\xi\rangle$ and
$\langle \gamma_U, \xi\rangle = \langle \gamma, U^{\tp}\xi\rangle + \langle \xi, \int c\,d\nu\rangle$, the
two terms in $\langle \xi, \int c\,d\nu\rangle$ cancel, and
\[
  \psi_{A^{\mathrm S}_U,\nu_U,\gamma_U}(\xi) = \psi_{A^{\mathrm S},\nu,\gamma}(U^{\tp}\xi).
\]
The transform of $U_*\mu$ at $\xi$ is $\widehat\mu(U^{\tp}\xi) = \exp\psi_{A^{\mathrm S},\nu,\gamma}(U^{\tp}\xi)$,
so $(A^{\mathrm S}_U, \nu_U, \gamma_U)$ is a generating triplet of $U_*\mu$, the only one by
Theorem~\ref{thm:cited-A39}(ii).

Every step is machine-checked. The formal statement, that $U_*\mu$ has this triplet, reads only
the form of the transform in Theorem~\ref{thm:cited-A39}(i) and no admitted fact. The
uniqueness in the last sentence is \ledger{A39}(ii).
\end{proof}

The next two theorems are [\sscite{sato1999levy}, Prop.~15.5, p.~93], with the sentence after it,
and [\sscite{sato1999levy}, Thm.~15.8, p.~94].

% shared with blueprint thm:cited-A43
\begin{theorem}[self-decomposable laws are infinitely divisible]
\label{thm:cited-A43}

If $\mu$ is self-decomposable, then $\mu$ is infinitely divisible and, for every $b > 1$, the
measure $\varrho_b$ of Definition~\ref{def:self-decomposable} is uniquely determined and infinitely
divisible. If $\mu$ is semi-selfdecomposable, then $\mu$ is infinitely divisible and $\varrho_b$ is
uniquely determined. So a self-decomposable law is semi-selfdecomposable with every $b > 1$ as a
span.
\end{theorem}

% shared with blueprint thm:cited-A44
\begin{theorem}[semi-selfdecomposability by the \Levy{} measure]
\label{thm:cited-A44}

Fix $b > 1$. Let $\mu$ be an infinitely divisible probability measure on $\RR^d$ with generating
triplet $(A^{\mathrm S}, \nu, \gamma)$. Then $\mu$ is semi-selfdecomposable with $b$ as a span if and only if
\[
  D_b\nu \ge \nu ,
\]
where $(D_b\nu)(B) = \nu(b^{-1}B)$ is the image of $\nu$ under $x \mapsto bx$. Semi-selfdecomposability
imposes no restriction on $A^{\mathrm S}$ and $\gamma$.
\end{theorem}

\begin{proof}
For a generating triplet $(A^{\mathrm S}, \nu, \gamma)$ write $\psi_{A^{\mathrm S},\nu,\gamma}(z)$ for the bracket of
Theorem~\ref{thm:cited-A39}(i), so that the law with this triplet has the transform
$\exp\psi_{A^{\mathrm S},\nu,\gamma}$. If $(A^{\mathrm S}_1, \nu_1, \gamma_1)$ and
$(A^{\mathrm S}_2, \nu_2, \gamma_2)$ are triplets, so is
$(A^{\mathrm S}_1 + A^{\mathrm S}_2, \nu_1 + \nu_2, \gamma_1 + \gamma_2)$: the sum of the matrices is symmetric and
nonnegative-definite, and $\nu_1 + \nu_2$ has no mass at $0$ and integrates $\min(|x|^2, 1)$. The
integrand of the bracket is integrable for $\nu_1$ and for $\nu_2$, so its integral against
$\nu_1 + \nu_2$ is the sum of the two integrals, and
\[
  \psi_{A^{\mathrm S}_1 + A^{\mathrm S}_2,\,\nu_1 + \nu_2,\,\gamma_1 + \gamma_2}
  = \psi_{A^{\mathrm S}_1,\nu_1,\gamma_1} + \psi_{A^{\mathrm S}_2,\nu_2,\gamma_2}.
\]

Let $\mu_b$ be the image of $\mu$ under $x \mapsto b^{-1}x$, so that
$\widehat\mu_b(z) = \widehat\mu(b^{-1}z)$. By Theorem~\ref{thm:cited-A42} with $n = d$ and $U = b^{-1}I$,
the law $\mu_b$ is infinitely divisible with a triplet $(b^{-2}A^{\mathrm S}, \nu_b, \gamma')$, where $\nu_b$
is the image of $\nu$ under $x \mapsto b^{-1}x$ (the restriction to $\RR^d \setminus \{0\}$ changes
nothing, as $\nu(\{0\}) = 0$) and $\gamma' \in \RR^d$. Since $\nu_b(B) = \nu(bB)$ and
$(D_b\nu)(B) = \nu(b^{-1}B)$, the inequality $D_b\nu \ge \nu$ on all Borel sets is the inequality
$\nu \ge \nu_b$, read at the set $bB$.

Let $\mu$ be semi-selfdecomposable with span $b$. By Definition~\ref{def:self-decomposable} its
residual $\varrho_b$ is infinitely divisible, so it has a triplet
$(A^{\mathrm S}_\varrho, \nu_\varrho, \gamma_\varrho)$ by
Theorem~\ref{thm:cited-A39}(i). Then
$\widehat\mu = \widehat\mu_b\,\widehat\varrho_b = \exp(\psi_{b^{-2}A^{\mathrm S},\nu_b,\gamma'} +
\psi_{A^{\mathrm S}_\varrho,\nu_\varrho,\gamma_\varrho})$, and by the first paragraph this is the transform given by the
summed triplet $(b^{-2}A^{\mathrm S} + A^{\mathrm S}_\varrho, \nu_b + \nu_\varrho, \gamma' + \gamma_\varrho)$. By
Theorem~\ref{thm:cited-A39}(ii) the summed triplet is $(A^{\mathrm S}, \nu, \gamma)$. So
$\nu = \nu_b + \nu_\varrho \ge \nu_b$, which is $D_b\nu \ge \nu$.

Conversely, let $\nu_b \le \nu$. The measure $\nu$ is $\sigma$-finite: it is finite on
$\{|x| > 1\}$ and on each shell $\{1/(n+1) < |x| \le 1/n\}$, $n \ge 1$, by the integrability of
$\min(|x|^2, 1)$, and $\nu(\{0\}) = 0$. From $\nu_b \le \nu$, every $\nu$-null set is
$\nu_b$-null, so by the Radon--Nikodym theorem $\nu_b = f\nu$ for a measurable $f \ge 0$, and
$f \le 1$ $\nu$-almost everywhere, again because $\nu_b \le \nu$. Put $\nu_\varrho := (1 - f)\nu$.
Then $\nu_b + \nu_\varrho = \nu$ and $\nu_\varrho \le \nu$, so $\nu_\varrho$ has no mass at $0$ and
integrates $\min(|x|^2, 1)$. The matrix
$(1 - b^{-2})A^{\mathrm S}$ is symmetric and nonnegative-definite because $b > 1$. By
Theorem~\ref{thm:cited-A39}(iii) there is an infinitely divisible law $\varrho_b$ with the triplet
$((1 - b^{-2})A^{\mathrm S}, \nu_\varrho, \gamma - \gamma')$. By the first paragraph
$\widehat\mu(b^{-1}z)\,\widehat\varrho_b(z) = \exp\psi_{A^{\mathrm S},\nu_b + \nu_\varrho,\gamma}(z) =
\widehat\mu(z)$, so $\mu$ is semi-selfdecomposable with span $b$ and residual $\varrho_b$.

The criterion involves $\nu$ alone, which is the last sentence of the statement. Every step is
machine-checked, and the cited facts are \ledger{A39}(i)--(iii) and \ledger{A42}.
\end{proof}

The polar form of a self-decomposable law is [\sscite{sato1999levy}, Thm.~15.10, p.~95].

% shared with blueprint thm:cited-A45
\begin{theorem}[the polar form of a self-decomposable law]
\label{thm:cited-A45}

Let $\mu$ be an infinitely divisible probability measure on $\RR^d$ with generating triplet
$(A^{\mathrm S}, \nu, \gamma)$, and let $S = \{\xi \in \RR^d : |\xi| = 1\}$. Then $\mu$ is self-decomposable if
and only if
\[
  \nu(B) = \int_S\lambda(d\xi)\int_0^\infty\mathbf{1}_B(r\xi)\,k_\xi(r)\,\frac{dr}{r}
\]
for every Borel set $B$, with a finite measure $\lambda$ on $S$ and a nonnegative function
$k_\xi(r)$, measurable in $\xi \in S$ and nonincreasing in $r > 0$. Self-decomposability imposes no
restriction on $A^{\mathrm S}$ and $\gamma$.
\end{theorem}

The next two theorems are on the support, [\sscite{sato1999levy}, Cor.~24.4, p.~149] and the
``only if'' of [\sscite{sato1999levy}, Cor.~24.6, p.~149]. The finiteness of $\nu$ in the second
follows from the support of $\nu$, as the proof below shows.

% shared with blueprint thm:cited-A46
\begin{theorem}[unbounded support]
\label{thm:cited-A46}

If $\mu$ is an infinitely divisible probability measure on $\RR^d$ and is not a point mass, then
the support of $\mu$ is unbounded.
\end{theorem}

% shared with blueprint thm:cited-A48
\begin{theorem}[infinitely divisible laws carried by $\ZZ$]
\label{thm:cited-A48}

Let $\mu$ be an infinitely divisible probability measure on $\RR$ with generating triplet
$(A^{\mathrm S}, \nu, \gamma)$. If the support of $\mu$ is contained in $\ZZ$, then $A^{\mathrm S} = 0$ and the
support of $\nu$ is contained in $\ZZ$. In particular $\nu(\RR) < \infty$.
\end{theorem}

\begin{proof}
The first sentence is cited. For the second: $\nu(\{0\}) = 0$ and $\nu$ is carried by $\ZZ$, so it
is carried by $\ZZ \setminus \{0\}$, where $|x| \ge 1$ and so $\min(|x|^2, 1) = 1$. Hence
$\nu(\RR) \le \int \min(|x|^2, 1)\,\nu(dx) < \infty$ by Theorem~\ref{thm:cited-A39}(i). Every step
is machine-checked, and the cited fact is \ledger{A48}, for the first sentence.
\end{proof}

Moments are read from the \Levy{} measure by the $g$-moment theorem [\sscite{sato1999levy},
Thm.~25.3, p.~159], with Definitions~25.1--25.2 on the same page. It is printed there for a
\Levy{} process and read here for its law at time $1$.

% shared with blueprint thm:cited-A49
\begin{theorem}[the $g$-moment theorem]
\label{thm:cited-A49}

A function $g$ on $\RR^d$ is \emph{submultiplicative} if it is nonnegative and there is a constant
$a > 0$ with $g(x + y) \le a\,g(x)g(y)$ for all $x, y \in \RR^d$, and \emph{locally bounded} if it is
bounded on every compact set. Let $g$ be a submultiplicative, locally bounded, measurable function on
$\RR^d$, and let $\mu$ be an infinitely divisible probability measure on $\RR^d$ with \Levy{} measure
$\nu$. Then $\int g\,d\mu < \infty$ if and only if $\int_{\{|x| > 1\}}g\,d\nu < \infty$.
\end{theorem}

The logarithmic moments of the grid pyramids use one function of Sato's list
[\sscite{sato1999levy}, Prop.~25.4(iii), p.~159], and the paper proves that it is
submultiplicative, with the constant $2$. For $g(x) = \log(|x| \vee e)$ and $u, v \ge e$ one has
$u + v \le uv$, so $|x + y| \vee e \le (|x| \vee e)(|y| \vee e)$. Hence
$g(x + y) \le g(x) + g(y) \le 2g(x)g(y)$, both values being at least $1$.

% shared with blueprint thm:cited-A50
\begin{theorem}[a logarithmic submultiplicative function]
\label{thm:cited-A50}

The function $g(x) = \log(|x| \vee e)$ on $\RR^d$ is submultiplicative.
\end{theorem}

\subsubsection{Bernstein functions}
\label{subsec:cited-bernstein}

A step of the \Matern{} pyramid in $\RR^d$ has a transform $e^{-f(|\omega|^2)}$, and
Section~\ref{ch:matern} says which class the function $f$ falls in
(Proposition~\ref{prop:matern-increments}). This subsection fixes the three classes it uses. A
function $f : (0,\infty) \to \RR$ is \emph{completely monotone} if it is smooth
and $(-1)^nf^{(n)} \ge 0$ for every $n \ge 0$. A \emph{Bernstein function} is a
function $f \ge 0$ on $(0,\infty)$ whose derivative is completely monotone. These are exactly the
functions
\[
  f(\lambda) = a + b\lambda + \int_{(0,\infty)}(1 - e^{-\lambda t})\,\eta(dt), \qquad a, b \ge 0,
  \quad \int_{(0,\infty)}\min(1, t)\,\eta(dt) < \infty,
\]
and the triplet $(a, b, \eta)$ is determined by $f$ [\sscite{schilling2012bernstein}, Thm.~3.2,
p.~21]. The measure $\eta$ is the \emph{\Levy{} measure} of $f$.

% shared with blueprint def:complete-bernstein
\begin{definition}[complete Bernstein functions]
\label{def:complete-bernstein}

A Bernstein function $f$ is a \emph{complete Bernstein function}
[\sscite{schilling2012bernstein}, Def.~6.1, p.~69] if its \Levy{} measure has a
completely monotone density $m$ with respect to Lebesgue measure:
\[
  f(\lambda) = a + b\lambda + \int_{(0,\infty)}(1 - e^{-\lambda t})\,m(t)\,dt .
\]
\end{definition}

% shared with blueprint def:thorin-bernstein
\begin{definition}[Thorin--Bernstein functions]
\label{def:thorin-bernstein}

A Bernstein function $f$ is a \emph{Thorin--Bernstein function}
[\sscite{schilling2012bernstein}, Def.~8.1, p.~109] if its \Levy{} measure has a density
$m$ such that $t \mapsto t\,m(t)$ is completely monotone, that is
\[
  f(\lambda) = a + b\lambda + \int_{(0,\infty)}(1 - e^{-\lambda t})\,m(t)\,dt,
\]
where $a, b \ge 0$, $\int_{(0,\infty)}\min(1, t)\,m(t)\,dt < \infty$ and $t \mapsto t\,m(t)$ is
completely monotone.
\end{definition}

A completely monotone function is a Laplace transform, by Bernstein's theorem
[\sscite{schilling2012bernstein}, Thm.~1.4, p.~3].

% shared with blueprint thm:cited-A55
\begin{theorem}[Bernstein's theorem]
\label{thm:cited-A55}

Let $f : (0,\infty) \to \RR$ be completely monotone. Then there is a measure $\mu$ on $[0,\infty)$
with
\[
  f(\lambda) = \int_{[0,\infty)}e^{-\lambda t}\,\mu(dt) \qquad\text{for every } \lambda > 0 .
\]
\end{theorem}

A continuous function $\psi$ on $\RR^d$ is \emph{negative definite} when $e^{-t\psi}$ is positive
definite for every $t > 0$. By [\sscite{schilling2012bernstein}, Thm.~13.14, p.~212], in the
direction proved by Example~13.13 on the same page, a Bernstein function of $|\omega|^2$ is negative
definite. The infinite divisibility of the law comes from subordination: Brownian motion run at the
random time of a subordinator, an increasing \Levy{} process [\sscite{sato1999levy}, Thm.~30.1,
pp.~197--198].

% shared with blueprint thm:cited-A57
\begin{theorem}[a Bernstein function of $|\omega|^2$]
\label{thm:cited-A57}

Let $f$ be a Bernstein function with $f(0{+}) = 0$, and set $f(0) := f(0{+}) = 0$. Then, for
every $d \ge 1$, the function
$\omega \mapsto f(|\omega|^2)$ is continuous and negative definite on $\RR^d$, and
$e^{-f(|\omega|^2)}$ is the transform of an infinitely divisible probability measure on $\RR^d$.
\end{theorem}

\subsubsection{\Polya{} frequency densities}
\label{subsec:cited-karlin}

The \Polya{} frequency densities of Section~\ref{subsec:cited-released} have reciprocal Laplace
transforms of one form, by Schoenberg's theorem [\sscite{karlin1968total}, Ch.~7, Thm.~3.2(a),
p.~345], with the class of entire functions of [\sscite{karlin1968total}, (2.2), p.~336]. The
identity of the next theorem relates the two-sided Laplace transform $L$ of the density to an
entire function $\psi$. Karlin states it on an open strip $\alpha < \operatorname{Re}s < \beta$
that contains the imaginary axis, on which $L$ exists [\sscite{karlin1968total}, Prop.~1.4,
p.~333]. The next theorem is narrower than the page: it states the identity on the imaginary axis
only, which is all the proofs read.

% shared with blueprint thm:cited-A24
\begin{theorem}[Schoenberg's representation]
\label{thm:cited-A24}

Let $\phi$ be a \Polya{} frequency density on $\RR$, and $L(s) := \int_{\RR}e^{-sx}\phi(x)\,dx$ its
two-sided Laplace transform. Then there are numbers $\gamma \ge 0$, $\delta \in \RR$ and real $a_j$
with $\sum_j a_j^2 < \infty$ and $\gamma + \sum_j a_j^2 > 0$ such that
\[
  L(i\omega)\,\psi(i\omega) = 1 \quad (\omega \in \RR), \qquad
  \psi(s) = e^{-\gamma s^2 + \delta s}\prod_j(1 + a_j s)\,e^{-a_j s},
\]
where $\psi$ is an entire function with $\psi(0) = 1$.
\end{theorem}

For a symmetric density the identity takes a real form, a consequence drawn here and not part of the
citation. The transform $L(i\omega) = \widehat\phi(\omega)$ is then real, nowhere $0$ by the identity, and
equal to $1$ at $0$, so it is positive. Taking moduli, with $|e^{i\delta\omega}| = 1$ and
$|(1 + ia_j\omega)e^{-ia_j\omega}| = (1 + a_j^2\omega^2)^{1/2}$, gives
\[
  -\log\widehat\phi(\omega) = \gamma\omega^2 + \tfrac12\sum_j\log(1 + a_j^2\omega^2) .
\]
The proof of
Proposition~\ref{prop:classical-pyramids-approximate-discrete-gaussian} in the supplement reads the
theorem this way.

\subsubsection{The modified Bessel function}
\label{subsec:cited-bessel}

The paper takes the modified Bessel function $K_\nu$ to be given by the integral of
[\sscite{dlmf2026}, 10.32.9], and cites the integral representation [\sscite{dlmf2026}, 10.32.10]
as the next theorem. The source states the representation for complex $z$ with
$|\mathrm{ph}\,z| < \pi/4$. The theorem is narrower, at real $z > 0$.

% shared with blueprint thm:cited-A51
\begin{theorem}[the modified Bessel function $K_\nu$]
\label{thm:cited-A51}

For real $\nu$ and real $z > 0$ let
$K_\nu(z) := \int_0^\infty e^{-z\cosh u}\cosh(\nu u)\,du$, the modified Bessel function of the
second kind of order $\nu$ (DLMF~10.32.9). Then for every real $\nu$ and every real $z > 0$
\[
  K_\nu(z) = \tfrac12\bigl(\tfrac12 z\bigr)^{\nu}\int_0^\infty
  \exp\Bigl(-t - \frac{z^2}{4t}\Bigr)\,\frac{dt}{t^{\nu+1}}
\]
(DLMF~10.32.10).
\end{theorem}
